\section{Documentation with Javadoc}\label{sec:javadoc}

\subsection{Comments in the code}

Every class, interface, enum and record, every constructor and every method carries a Javadoc comment, private
members and the handlers the FXML views call included. The comments follow the conventions of the documentation
comment specification~\cite{doc-comment-spec}:
\begin{itemize}
  \item the first sentence says what the element is or does, and becomes its summary in the generated pages;
  \item each parameter is described with \texttt{@param}, and the result of every method that returns one with
        \texttt{@return};
  \item an exception the method throws on purpose is described with \texttt{@throws}, saying when;
  \item every type names its two authors, one \texttt{@author} tag each: \texttt{@author Hamza Rasool} and
        \texttt{@author Abdullah Khan}.
\end{itemize}
Listing~\ref{lst:javadoc} shows the class comment of the task service and the comment of its \texttt{create} method,
as they are in the source.

\source[linerange={15-29,91-108}, gobble=0]{src/main/java/io/mcs/taskmanager/services/TaskService.java}
  {Javadoc comments in \texttt{TaskService}: the class and its \texttt{create} method}{lst:javadoc}
% Check after any change to TaskService: lines 15-29 are the class comment and 91-108 create.

\subsection{Generating the HTML}

The HTML documentation is generated by the Javadoc tool~\cite{javadoc-tool} through the Maven Javadoc Plugin, with one
command from the project folder:
\begin{center}
  \texttt{./mvnw javadoc:javadoc}
\end{center}
The plugin's configuration in \texttt{pom.xml} (Listing~\ref{lst:javadoc-plugin}) turns on every check doclint has
and makes any warning fail the build rather than scroll past, includes private members, shows the \texttt{@author}
tags, and leaves the date out of the pages, so that generating them again changes only what the code changed. The
pages are written into the repository, under \path{docs/javadoc/apidocs/}, and committed with the code, so they are
part of the submission.

\source[firstline=106, lastline=126, gobble=3, tabsize=4, language=XML]{pom.xml}
  {The Maven Javadoc Plugin's configuration in \texttt{pom.xml}}{lst:javadoc-plugin}

Because a single warning would fail it, a run that ends in BUILD SUCCESS means that every comment passed every check.
The last run before this report was built printed no warning at all:
\begin{lstlisting}[style=terminal, literate={-}{{-\allowbreak}}1]
$ ./mvnw javadoc:javadoc
...
[INFO] --- javadoc:3.12.0:javadoc (default-cli) @ taskmanager ---
[INFO] ------------------------------------------------------------------------
[INFO] BUILD SUCCESS
[INFO] ------------------------------------------------------------------------
[INFO] Total time:  3.302 s
[INFO] Finished at: 2026-11-22T22:50:32+05:00
\end{lstlisting}

\screenshotwide{javadoc-01-generate}
  {Generating the documentation with \texttt{./mvnw javadoc:javadoc}, ending in BUILD SUCCESS.}
  {Terminal: ./mvnw javadoc:javadoc ending in BUILD SUCCESS}

\subsection{Reviewing it in a browser}

The generated pages were opened from \path{docs/javadoc/apidocs/index.html} in a web browser and reviewed there. The
overview lists the seven packages, \texttt{io.mcs.taskmanager} and its \texttt{entities}, \texttt{repositories},
\texttt{services}, \texttt{shared}, \texttt{ui} and \texttt{ui.kit}, each with the summary from its
\texttt{package-info.java} (Fig.~\ref{fig:javadoc-02-index}). Each class page lists its fields, constructors and
methods, private ones included, with their parameters, return values and exceptions, and names the class's authors
(Fig.~\ref{fig:javadoc-03-class}). The documentation is 71 HTML pages in all, with an index and a search.
% VERIFY in phase 3: the screenshots show this.

\screenshotwide{javadoc-02-index}
  {The generated documentation's overview page, \texttt{docs/javadoc/apidocs/index.html}, in a browser.}
  {docs/javadoc/apidocs/index.html in a browser}

\screenshotwide{javadoc-03-class}
  {The task service's page: method details with their parameters and return values.}
  {The task service's page in the browser, method details with parameters and returns}
